\documentclass[12pt]{article}
\usepackage[utf8]{inputenc}
\usepackage[T1]{fontenc}
\usepackage{hyperref}
\title{Tarpia Language Resources for K-12 Curriculum Development}
\author{Compiled by Claude Sonnet 5 (claude-sonnet-5) for Andr\'as Kornai}
\date{}
\begin{document}
\maketitle

\section{Introduction}
Tarpia (also Tarfia or Sufrai; ISO 639-3 tpf) is an Austronesian language of the Sarmi branch of Oceanic, spoken on the eastern north coast of Papua province, Indonesia. Sources give 630 native speakers (2006) or 564 (data dated 1978), and call it endangered. Joshua Project places the people in Tarfia village, Demta subdistrict, Jayapura Regency, and classifies the language differently (as Lakes Plain), which conflicts with the other sources. Very little exists, and what follows is short.

\section{Primary Educational Resources}
None found in search.

\section{Linguistic and Reference Materials}
\begin{itemize}
\item Wikipedia, Tarpia language: \url{https://en.wikipedia.org/wiki/Tarpia_language}
\item Wikipedia, Sarmi--Jayapura languages: \url{https://en.wikipedia.org/wiki/Sarmi%E2%80%93Jayapura_languages}
\item Endangered Languages Project entry: \url{https://endangeredlanguages.com/lang/2790}
\item Glottolog reference record for Kim (2006), a SIL survey report covering Tarpia along with Bgu and Kaptiau: \url{https://glottolog.org/resource/reference/id/60643}
\end{itemize}

\section{Religious and Community Resources}
\begin{itemize}
\item Joshua Project profile (notes no complete Bible translation in Tarpia and that younger people increasingly use Indonesian): \url{https://joshuaproject.net/people_groups/15275}
\end{itemize}

\section{Digital and Audio Resources}
None found in search.

\section{Curriculum Development Context}
The language is not known to be taught in schools, and no teaching materials, dictionary or grammar were found. With speakers mostly older adults, documentation and intergenerational transmission would come before curriculum work. The SIL survey report by Kim (2006) is the most useful lead for further research.
\end{document}